% Source fingerprint: b498fd314c29f9c0aadfeeb4aca942a02d5a32b6174e8cf3848275fbd1ed6b88
% T13: raw TeX cells preserved; no numeric highlighting.
\begin{table}[htbp]
\centering\small\renewcommand{\arraystretch}{1.18}\setlength{\tabcolsep}{4pt}
\caption[高阶 Sobolev 嵌入速查]{高阶 Sobolev 嵌入速查。本表补充前面的一阶指数表。函数空间的范数与代表约定均沿用本节，不把内部正则性写成跨边界结论。}
\label{b1:tab:visual-t13}
\begin{tabularx}{\linewidth}{@{}>{\raggedright\arraybackslash}p{3.1cm}>{\raggedright\arraybackslash}p{4.7cm}>{\raggedright\arraybackslash}X@{}}
\toprule
参数区间 & 本节得到的目标 & 适用说明 \\
\midrule
\(kp<d\)，\(1\le p<\infty\) & \(L^q\)，\(q=dp/(d-kp)\) & 逐次一阶嵌入；分母必须严格为正 \\
\(kp=d\) & 不能把上行公式代入为无穷指数 & 需按相应临界理论处理，不由形式代入宣称有界连续 \\
\(k-d/p=m+\alpha\)，\(0<\alpha<1\) & 具有 \(C^{m,\alpha}\) 代表 & 有限 p；余量非整数，经典导数的相容性由正文证明 \\
\(j=k-d/p\in\mathbb N_{>0}\) & \(C^{j-1,\alpha}\)，任意 \(0<\alpha<1\) & 本节该命题要求 \(1<p<\infty\)；不擅取 \(\alpha=1\) \\
一般区域 & 内部结论通过局部化取得；全域结论需相应高阶延拓 & 一阶延拓不能直接替代 \(k\) 阶延拓假设 \\
\bottomrule
\end{tabularx}
\end{table}
